\section{Harmless quotient owners in the full invariant lattice}
\label{sec:harmless-exact-owners}

% C6: source-preserved, no change of theorem signature or proof.  This
% general full-invariant-lattice statement is used both in the
% single-minimal-infinite-cut reduction and in the maximal exact core
% construction; it requires only the earlier real-hull theorem.
\begin{lemma}[Harmless owners contain all minimal directions]
\label{lem:acf-harmless-owners}
In any standing infinite source, naming finitely many quotient
classes of an invariant equivalence $H$ does not increase
\emph{real} algebraic closure if and only if $H$ contains
all inclusion-minimal infinite equivalences.  The same holds
simultaneously for finitely many such sorts.
\end{lemma}
\begin{proof}
We prove both directions, distinguishing
the FULL invariant-equivalence lattice
from its exact sublattice.

First assume \(H\) contains ALL
inclusion-minimal INFINITE invariant
equivalences. Fix finite actual real
\(A\), finitely many actual named
\(H\)-classes \(C_1,\ldots,C_m\)
(or finitely many such classes in
several owner sorts, each owner kernel
containing every minimal infinite
equivalence), and a real point
\(x\notin\operatorname{acl}_M(A)\).
We show
\[
     x\notin
     \operatorname{acl}_M(A,C_1,\ldots,C_m).
\]

For each \(a\in A\) put
\(p_a=\col(x,a)\) (when \(x=a\),
which cannot happen here, the
equality relation would be used),
and define the ACTUAL invariant
equivalence
\[
       J_A=\bigcap_{a\in A}\Gamma(p_a).
\]
The class of \(x\) in this
equivalence is
\[
       J_A[x]=
       \bigcap_{a\in A}
                   [a]_{\Gamma(p_a)},
\]
because \(x\Gamma(p_a)a\).
If it were finite it would be
a finite \(A\)-definable set
containing \(x\), contradicting
\(x\notin\operatorname{acl}(A)\).
Thus \(J_A\) has infinite classes.
The invariant-equivalence lattice
is finite (unions of finitely many
old atomic orbitals), so it contains
some inclusion-minimal infinite
equivalence \(E\le J_A\).
We fix ONE such \(E\) for ALL
owner classes to be named.

For any named owner class \(C_j\)
of a kernel \(H_j\ge E\), choose
an \(E\)-class \(D_j\subseteq C_j\).
It is infinite. Let \(F=\Sigma(\Id)\)
be the greatest finite-class
invariant equivalence of \(M\);
all its classes are finite.
Call an actual old atomic colour
\(p\) BAD (for this \(E\)) when
\(E\cap(p;p)\subseteq F\).
For such \(p\), any two points
\(b,b'\in p[x]\cap D_j\)
are related both by \(E\) and
by \(p;p\), hence by \(F\).
Therefore this entire intersection,
when nonempty, lies in ONE
finite \(F\)-class. There are
only finitely many old atomic
colours, so the union of these
bad neighbour sets meets
\(D_j\) in finitely many points.
Choose an actual real representative
\(b_j\in D_j\) OUTSIDE that finite
union AND outside the singleton
\(\{x\}\). This extra exclusion
is essential: \(x\) might itself
lie in \(D_j\), and the subsequent
\(p_j=\col(x,b_j)\) must be a
genuine DIVERSITY atomic colour,
not logical equality. Since
\(D_j\) is infinite and only
finitely many points are
excluded, such \(b_j\) exists.
Its genuine atomic diversity
colour \(p_j=\col(x,b_j)\)
is therefore not bad.

For this good \(p_j\),
\(E\cap(p_j;p_j)\nsubseteq F\).
The relation \(E\cap\Gamma(p_j)\)
is an invariant EQUIVALENCE,
contains \(E\cap(p_j;p_j)\)
and is therefore NOT a
finite-class equivalence, since
every finite-class invariant
equivalence lies below \(F\).
By minimality of the INFINITE
\(E\), this intersection must
be all of \(E\). Consequently
\[
       E\le\Gamma(\col(x,b_j))
                    \quad(j=1,\ldots,m).
\]
Since \(E\le J_A\) as well,
\[
 E[x]\subseteq
  \bigcap_{a\in A}[a]_{\Gamma(\col(x,a))}
  \cap
  \bigcap_{j=1}^m[b_j]_{\Gamma(\col(x,b_j))}.
\tag{HO1}\label{eq:acf-harmless-hull}
\]
This intersection is INFINITE.
If \(x\) were algebraic over
\(A\) and the NAMED imaginary
owner classes \(C_1,\ldots,C_m\),
it would also be algebraic over
\(A\cup\{b_1,\ldots,b_m\}\),
because each \(C_j\) is definable
from its chosen representative
\(b_j\). The canonical real-hull
theorem \cref{thm:rh-hull}
would then force the displayed
intersection to be FINITE,
contrary to \eqref{eq:acf-harmless-hull}.
This proves harmlessness of
ANY finite collection of such
owner classes simultaneously.

Conversely suppose \(H\) omits
some minimal infinite \(E\):
\(E\nleq H\). The invariant
intersection \(E\cap H\)
is a proper equivalence below
the minimal infinite \(E\),
hence has FINITE real classes
and lies below \(F\).
Choose actual points \(x,u\)
with \(uEx\) but \(u\not F x\),
possible because \(E\) has
infinite classes while \(F\)
has finite classes. Then
\[
       E[u]\cap H[x]
                =(E\cap H)[x]
\]
is a finite nonempty set,
definable over the real point
\(u\) and the imaginary class
\(H[x]\), and contains \(x\).
Thus \(x\in\operatorname{acl}
(u,H[x])\). On the other hand
\(\operatorname{acl}(u)=F[u]\).
For clarity, this singleton
identity has an elementary
finite-orbital proof. Every
\(F\)-class is finite and
\(F[u]\) is \(u\)-definable,
so \(F[u]\subseteq\operatorname{acl}(u)\).
Conversely, if an atomic
suborbit \(p[u]\) is finite,
the symmetric orbital
\(p\)-graph is LOCALLY FINITE,
with all vertices having
the same finite degree.
Its connectedness equivalence
\(\Gamma(p)\) is invariant
and thus a union of FINITELY
many actual old atomic
orbitals. For each such
orbital \(q\subseteq\Gamma(p)\),
choose a single realised
\(q\)-pair and a finite
\(p\)-walk between its
endpoints. Homogeneity on
the OLD pair orbital \(q\)
shows the same walk LENGTH
works for every \(q\)-pair.
Taking the maximum over
finitely many \(q\)'s gives
a UNIFORM finite diameter
for every connected
\(p\)-component. Local
finiteness then makes every
component FINITE. Thus
\(\Gamma(p)\) is a finite-class
invariant equivalence and
hence lies below the
greatest such \(F\).
Every real algebraic
single-point suborbit is
a union of finitely many
finite atomic suborbits,
so it lies in \(F[u]\).
Therefore
\(\operatorname{acl}(u)=F[u]\),
and \(x\notin\operatorname{acl}(u)\).
Naming this one \(H\)-class
increases real algebraicity,
proving necessity.
\end{proof}
